\begin{center}
    {\LARGE \textbf{LỜI GIẢI ĐỀ LUYỆN 04}}\\[6pt]
    {\large (Dành cho đội tuyển quốc gia)}
\end{center}

\section*{Bài 1: Chia và nối mảng}
% https://www.codechef.com/problems/SPLITADD

Quá trình được mô tả trong đề bài là đệ quy, vì vậy ta cũng cố gắng sử dụng tư duy đệ quy để giải.

\subsection*{Subtask 1: Tính một phần tử duy nhất}

Giả sử ta muốn tìm đúng một phần tử, tức là $K = L = R$. Gọi:
\[
ans(N, K)
\]
là giá trị tại vị trí $K$ trong mảng kết quả có độ dài $N$.

Hàm $f(A)$ chia mảng thành hai nửa: các phần tử ở vị trí lẻ và các phần tử ở vị trí chẵn, sau đó đệ quy trên từng nửa.

Do đó, có hai trường hợp:

\paragraph{Trường hợp 1: $K \le \frac{N}{2}$}

Ta đang tìm phần tử thứ $K$ trong mảng:
\[
f([1,3,5,\dots,N-1]).
\]

Mảng này gồm toàn số lẻ, và phần tử thứ $i$ có dạng:
\[
2i - 1.
\]

Vì vậy, ta có thể tìm phần tử thứ $K$ của:
\[
f([1,2,3,\dots,\frac{N}{2}])
\]
sau đó nhân đôi và trừ đi $1$.

\[
ans(N,K) = 2 \cdot ans\!\left(\frac{N}{2}, K\right) - 1.
\]

\paragraph{Trường hợp 2: $K > \frac{N}{2}$}

Ta tìm phần tử thứ:
\[
K - \frac{N}{2}
\]
của:
\[
f([2,4,6,\dots,N]).
\]

Mảng này chính là mảng $[1,2,\dots,\frac{N}{2}]$ nhưng nhân đôi mỗi giá trị.

Do đó:
\[
ans(N,K) = 2 \cdot ans\!\left(\frac{N}{2}, K - \frac{N}{2}\right).
\]

\paragraph{Trường hợp cơ sở:}
\[
ans(1,1) = 1.
\]

\paragraph{Độ phức tạp:}
Tại mỗi bước ta chia đôi $N$, nên độ phức tạp là:
\[
\mathcal{O}(\log N).
\]

\subsection*{Subtask 2: Tính tổng một đoạn $[L,R]$}

Ta mở rộng hàm thành:
\[
ans(N,L,R)
\]
là tổng giá trị các phần tử từ vị trí $L$ đến $R$ của mảng kết quả.

Đặt $H = \frac{N}{2}$. Có ba trường hợp:

\paragraph{1. Toàn bộ đoạn nằm trong nửa trái ($R \le H$):}
\[
ans(N,L,R)
= 2 \cdot ans(H,L,R) - (R-L+1).
\]

\paragraph{2. Toàn bộ đoạn nằm trong nửa phải ($L > H$):}
\[
ans(N,L,R)
= 2\cdot ans(H, L-H, R-H).
\]

\paragraph{3. Đoạn giao giữa hai nửa ($L \le H < R$):}

Ta tách thành hai đoạn:
\[
ans(N,L,R) = ans(N,L,H) + ans(N,H+1,R).
\]

\subsection*{Tối ưu quan trọng}

Nếu đoạn truy vấn bao phủ toàn bộ mảng:
\[
ans(N,1,N) = \frac{N(N+1)}{2},
\]
ta trả về ngay mà không tiếp tục đệ quy.

Điều này giảm số nhánh đệ quy xuống còn $\mathcal{O}(\log N)$.

\subsection*{Độ phức tạp cuối cùng}
\[
\mathcal{O}(\log N)
\]
cho mỗi test.

Đừng quên lấy modulo $10^9+7$ và tránh tràn số khi tính $\frac{N(N+1)}{2}$ bằng cách lấy $N \bmod mod$ trước khi nhân.

\newpage

\section*{Bài 2. Rắn săn mồi}
% https://www.codechef.com/problems/SNAGR

\subsection*{Giải một truy vấn đơn lẻ}

Giả sử ta có sẵn một ma trận Boolean $A$ kích thước $N \times N$:

\[
A_{i,j} = 
\begin{cases}
1 & \text{nếu ô $(i,j)$ đang bị một con rắn chiếm giữ}, \\
0 & \text{ngược lại}.
\end{cases}
\]

Nhiệm vụ là tìm đường đi có chi phí nhỏ nhất từ $(1,1)$ đến $(N,N)$, trong đó:

\[
\text{cost(path)} = \sum_{\text{cells } (i,j) \in \text{path}} A_{i,j}
\]

và ta được phép đi 4 hướng: lên, xuống, trái, phải.

Đây tương ứng với một đồ thị có $N^2$ đỉnh và $\le 4N^2$ cạnh.  
Dijkstra có thể giải trong thời gian:

\[
\mathcal{O}(N^2 \log N),
\]

nhưng vẫn quá chậm.

\subsection*{Sử dụng 0--1 BFS}

Trong bài này, mỗi cạnh có trọng số chỉ là $0$ hoặc $1$ (ô trống hoặc ô nguy hiểm).  
Do đó, ta có thể dùng thuật toán \textbf{0--1 BFS}, chạy trong:

\[
\mathcal{O}(N^2)
\]

vì chỉ dùng deque và mỗi cạnh được xử lý đúng 1 lần.

Như vậy, ta có thể tìm độ nguy hiểm nhỏ nhất tại một thời điểm $T$ trong thời gian $\mathcal{O}(N^2)$.

\subsection*{Xử lý nhiều truy vấn}

Có tối đa $Q = 2 \cdot 10^5$ truy vấn, nên ta không thể chạy $\mathcal{O}(N^2)$ cho từng truy vấn.

Quan sát quan trọng:

\[
\text{Từ thời điểm } T = N \text{ trở đi, trạng thái lưới không thay đổi nữa.}
\]

Lý do: mỗi con rắn sẽ chạm vào biên sau nhiều nhất $N-1$ bước.

Suy ra:

- Ta chỉ cần tính đáp án cho $T = 0, 1, \ldots, N$.
- Mỗi giá trị $T$ tốn $\mathcal{O}(N^2)$.
- Tổng thời gian tiền xử lý:

\[
\mathcal{O}(N^3)
\]

Sau đó, mỗi truy vấn trả lời trong:

\[
\mathcal{O}(1)
\]

bằng cách lấy kết quả tiền tính.

\subsection*{Cách tính ma trận $A$ tại các thời điểm}

Cần sinh ra ma trận $A^{(T)}$ cho tất cả $T = 0 \ldots N$.

\subsubsection*{Cách đơn giản nhưng chậm}

Với mỗi ô chứa đầu rắn, duỗi nó $T-1$ bước theo hướng đã cho.  
Cách này tốn tối đa:

\[
\mathcal{O}(N^3)
\]

cho một $T$, và nếu làm cho mọi $T$ thì là $\mathcal{O}(N^4)$ — quá chậm.

\subsubsection*{Tối ưu hoá bằng mô phỏng gia tăng}

Khi tăng thời gian từ $T$ lên $T+1$:

- Mỗi con rắn chỉ di chuyển thêm đúng 1 ô.
- Ta chỉ cần cập nhật những ô mới bị chiếm.

Ta lưu:

- ma trận $A$ hiện tại,
- danh sách các đầu rắn và hướng di chuyển.

Với tối đa $N^2$ con rắn và $N$ bước thời gian:

\[
\mathcal{O}(N^3)
\]

\subsubsection*{Tối ưu hoá bằng tính toán thời điểm chiếm ô}

Một cách khác:

- Với mỗi ô $(i,j)$, tìm thời điểm đầu tiên nó bị rắn chạm tới.
- Ta xét tất cả rắn, duỗi theo hướng của nó cho đến khi chạm tường.

Lặp toàn bộ mất:

\[
\mathcal{O}(N^2)
\]

Sau đó, tại truy vấn $T$:

\[
A_{i,j}^{(T)} = 1 \iff \text{time\_first\_occupied}(i,j) < T
\]

Việc này rất nhanh và tiện lợi.

\subsection*{Độ phức tạp}

\[
\boxed{
\mathcal{O}(N^3 + Q)
}
\]

cho mỗi test.

\newpage

\section*{Bài 3: Ice-cream}

\subsection*{Tìm tập con có XOR lớn nhất}

Giả sử ta có một dãy số và cần chọn một tập con sao cho giá trị XOR của tập con đó đạt giá trị lớn nhất có thể.

Kí hiệu $L(S)$ là tập chứa tất cả các giá trị có thể thu được bằng cách XOR các phần tử thuộc mọi tập con của $S$, kể cả tập rỗng.  
Ví dụ:
\[
L(\{1,5\}) = \{0, 1, 4, 5\}.
\]
Mục tiêu của chúng ta là tìm phần tử lớn nhất của $L(S)$.

\subsection*{Bổ đề 1}
Cho $S$ là tập các số, và $S'$ là tập mới nhận được bằng cách lấy hai phần tử khác nhau $A, B \in S$, rồi thay chúng bằng $A$ và $A \oplus B$. Khi đó:
\[
L(S) = L(S').
\]

\paragraph{Ví dụ.}  
Nếu $S=\{4,6,7\}$ thì một khả năng của $S'$ là:
\[
\{4, 4\oplus 6, 7\}.
\]

\paragraph{Chứng minh.}  
Lấy $X \in L(S)$ bất kỳ. Ta có thể viết:
\[
X = X_1 \oplus X_2 \oplus \cdots \oplus X_k,
\]
trong đó mỗi $X_i$ là một phần tử của $S$.

Xét hai trường hợp:
\begin{enumerate}
    \item Nếu $B$ không xuất hiện trong dãy $\{X_i\}$ thì rõ ràng $X \in L(S')$ vì tất cả các $X_i$ đều nằm trong $S'$.
    \item Nếu $B$ có trong dãy $\{X_i\}$, ta thay $B$ bằng hai phần tử $A$ và $A\oplus B$. Khi đó:
    \[
    B = A \oplus (A \oplus B),
    \]
    nên biểu thức XOR của $X$ không thay đổi. Do đó $X$ vẫn thuộc $L(S')$.
\end{enumerate}

Vậy $L(S)=L(S')$.

\subsection*{Bổ đề 2}

Nếu $S$ chứa số $0$ và $S'$ là tập nhận được từ $S$ bằng cách loại bỏ $0$, thì:
\[
L(S)=L(S').
\]

\paragraph{Chứng minh.}  
Vì $0$ là phần tử trung hoà của XOR, nên thêm hoặc bỏ $0$ không làm thay đổi bất kỳ giá trị XOR nào.

\subsection*{Ý tưởng chính}

Từ hai bổ đề trên, ta có thể tự do biến đổi dãy số, miễn là:
\begin{itemize}
    \item thay một phần tử bởi XOR của nó với phần tử khác,
    \item loại bỏ các số $0$ phát sinh.
\end{itemize}

Ta muốn biến đổi dãy sao cho không có hai số nào có cùng vị trí bit cao nhất (MSB).  
Ta duyệt từng phần tử $C$ trong dãy theo thứ tự:

\begin{enumerate}
    \item Với mỗi phần tử trong tập cơ sở, nếu phần tử đó có cùng MSB với $C$, ta thay:
    \[
    C \leftarrow C \oplus V.
    \]
    \item Lặp lại cho đến khi:
    \begin{itemize}
        \item $C = 0$ (khi đó bỏ qua theo Bổ đề 2), hoặc
        \item không còn phần tử nào trong cơ sở trùng MSB với $C$, khi đó ta thêm $C$ vào cơ sở.
    \end{itemize}
\end{enumerate}

Vì mỗi lần thêm số mới sẽ tạo ra một MSB mới, mà các số đều $\le 10^9$ nên dãy cơ sở có tối đa 30 phần tử.

\subsection*{Tìm XOR lớn nhất}

Sau khi tạo được một cơ sở tuyến tính có dạng “mỗi số có MSB khác nhau”, ta có thể tìm giá trị XOR lớn nhất bằng cách tham lam:

\begin{enumerate}
    \item Đặt \texttt{ANS = 0}.
    \item Duyệt các chỉ số bit từ lớn xuống nhỏ (29 về 0).
    \item Nếu tồn tại số $X$ trong cơ sở có MSB tại vị trí đó và:
    \[
    (ANS \oplus X) > ANS,
    \]
    thì ta cập nhật:
    \[
    ANS \leftarrow ANS \oplus X.
    \]
\end{enumerate}

Kết thúc quá trình, \texttt{ANS} chính là XOR lớn nhất.

\subsection*{Cài đặt mẫu}

\begin{verbatim}
struct Gauss{
    int arr[31];
    int len;
    Gauss(){ len=0; }

    void add(int x){
        for(int i=0;i<len;i++){
            if((arr[i]^x) < x)
                x = arr[i] ^ x;
        }
        if(x != 0){
            arr[len++] = x;
            for(int i=len-1;i>0;i--){
                if(arr[i] > arr[i-1])
                    swap(arr[i], arr[i-1]);
            }
        }
    }

    int getMax(){
        int ret = 0;
        for(int i=0;i<len;i++){
            if((ret ^ arr[i]) > ret)
                ret ^= arr[i];
        }
        return ret;
    }
};
\end{verbatim}

\subsection*{Hỗ trợ xoá phần tử cuối cùng}

Ta có thể lưu mỗi trạng thái của cơ sở vào một stack.  
Mỗi lần thêm phần tử mới, ta sao chép cơ sở hiện tại, thêm phần tử, rồi đẩy vào stack.  
Để xoá phần tử thêm gần nhất, ta chỉ cần pop khỏi stack.

\subsection*{Hỗ trợ xoá phần tử sớm nhất}

Ta có thể xây dựng \href{https://cp-algorithms.com/data_structures/stack_queue_modification.html#queue-modification-method-3}{một queue bằng hai stack} như sau:

\begin{itemize}
    \item Stack \texttt{forward} lưu các phần tử mới thêm.
    \item Stack \texttt{backward} lưu các phần tử theo thứ tự bị đảo ngược.
    \item Khi cần xoá phần tử lâu nhất, nếu \texttt{backward} rỗng thì ta chuyển toàn bộ phần tử từ \texttt{forward} sang \texttt{backward}.
\end{itemize}

Trong lời giải, mỗi phần tử của queue chính là một cơ sở có tối đa 30 phần tử, nên chi phí gộp hai cơ sở là hằng số nhỏ.

\subsection*{Tổng quan lời giải đầy đủ}

Ta duy trì hai con trỏ tạo thành một đoạn trượt.  
Mỗi lần tăng con trỏ phải, ta thêm phần tử.  
Mỗi lần tăng con trỏ trái, ta xoá phần tử sớm nhất.  
Mỗi khi có đoạn hợp lệ (tổng chi phí $\le K$), ta lấy hai cơ sở từ hai stack, gộp chúng lại, rồi tính XOR lớn nhất.

Độ phức tạp:
\[
O(N \log^2(10^9)) \approx O(N \cdot 30 \cdot 30),
\]
hoàn toàn đủ nhanh.
